% =============================================================
% TikZ Şekil Makroları - Bilgisayar Mimarisi Kitabı
% Oğuz Ergin
% =============================================================
% Bu dosya tüm bölümlerde kullanılan ortak TikZ şekil
% makrolarını ve stil tanımlarını içerir.
% =============================================================

% -----------------------------------------------
% RENKLER
% -----------------------------------------------
\definecolor{sinyalmavi}{RGB}{70, 130, 180}
\definecolor{sinyalkirmizi}{RGB}{200, 60, 60}
\definecolor{sinyalyesil}{RGB}{60, 160, 60}
\definecolor{blokgri}{RGB}{240, 240, 240}
\definecolor{blokmavi}{RGB}{200, 220, 240}
\definecolor{blokkirmizi}{RGB}{255, 220, 220}
\definecolor{blokyesil}{RGB}{220, 245, 220}
\definecolor{bloksari}{RGB}{255, 250, 205}
\definecolor{blokturuncu}{RGB}{255, 225, 180}
\definecolor{vurgumavi}{RGB}{0, 100, 200}
\definecolor{vurgukirmizi}{RGB}{200, 0, 0}
\definecolor{vurguyesil}{RGB}{0, 140, 60}
\definecolor{vurguturuncu}{RGB}{200, 120, 0}

% -----------------------------------------------
% RENK VARYANTLARI (açık tonlar)
% -----------------------------------------------
\colorlet{blokmavi-acik}{blokmavi!30}
\colorlet{blokkirmizi-acik}{blokkirmizi!30}
\colorlet{blokyesil-acik}{blokyesil!30}
\colorlet{bloksari-acik}{bloksari!30}
\colorlet{blokgri-acik}{blokgri!30}

% -----------------------------------------------
% BORU HATTI AŞAMA RENKLERİ
% -----------------------------------------------
\definecolor{asamagetir}{RGB}{200, 220, 240}    % G: Getir — mavi
\definecolor{asamacoz}{RGB}{220, 245, 220}       % Ç: Çöz — yeşil
\definecolor{asamayurut}{RGB}{255, 235, 200}     % U: Yürüt — turuncu
\definecolor{asamabellek}{RGB}{255, 220, 220}    % B: Bellek — kırmızı
\definecolor{asamayaz}{RGB}{230, 220, 250}       % Y: Yaz — mor

% -----------------------------------------------
% GENEL TikZ STİLLERİ
% -----------------------------------------------

% Blok şeması stilleri
\tikzset{
    % Temel blok stili
    blok/.style={
        draw, rectangle, rounded corners=2pt,
        minimum width=2cm, minimum height=1cm,
        fill=blokmavi, font=\small,
        align=center
    },
    % Vurgulu blok
    vurgublok/.style={
        blok, fill=blokkirmizi, draw=vurgukirmizi, thick
    },
    % Bellek bloku
    bellekblok/.style={
        draw, rectangle, minimum width=2.5cm, minimum height=1.5cm,
        fill=blokyesil, font=\small, align=center
    },
    % Yazmaç bloku
    yazmacblok/.style={
        draw, rectangle, minimum width=2cm, minimum height=0.8cm,
        fill=bloksari, font=\small, align=center
    },
    % Ok stilleri
    okstil/.style={
        ->, >=Stealth, thick
    },
    ciftok/.style={
        <->, >=Stealth, thick
    },
    veristil/.style={
        ->, >=Stealth, thick, sinyalmavi
    },
    denetimstil/.style={
        ->, >=Stealth, dashed, sinyalkirmizi
    },
}

% -----------------------------------------------
% KARNO HARİTASI MAKROSU
% -----------------------------------------------

% 2 değişkenli Karno haritası
% #1: değişken isimleri (örn. "A" ve "B")
% #2: hücre değerleri (4 adet, sırasıyla 00,01,10,11)
\newcommand{\karnoiki}[4][]{%
    % #2=değişken1, #3=değişken2, #4=değerler (virgülle ayrılmış 4 değer)
    \begin{tikzpicture}[scale=0.8]
        % Grid
        \draw (0,0) grid (2,2);
        % Değişken etiketleri
        \node[above] at (0.5, 2) {\footnotesize 0};
        \node[above] at (1.5, 2) {\footnotesize 1};
        \node[left] at (0, 1.5) {\footnotesize 0};
        \node[left] at (0, 0.5) {\footnotesize 1};
        \node[above left] at (0, 2) {\footnotesize #2\textbackslash #3};
        % Hücre değerleri
        \foreach \val [count=\i from 0] in {#4} {
            \pgfmathtruncatemacro{\col}{mod(\i,2)}
            \pgfmathtruncatemacro{\row}{1 - int(\i/2)}
            \node at (\col+0.5, \row+0.5) {\val};
        }
        #1
    \end{tikzpicture}
}

% 3 değişkenli Karno haritası
% #1: ek TikZ kodu (gruplamalar vb.)
% #2: satır değişkeni, #3: sütun değişkenleri
% #4: 8 hücre değeri (AB\C sıralaması: 00/0, 00/1, 01/0, 01/1, 11/0, 11/1, 10/0, 10/1)
\newcommand{\karnouc}[4][]{%
    \begin{tikzpicture}[scale=0.8]
        \draw (0,0) grid (4,2);
        \node[above] at (0.5, 2) {\footnotesize 00};
        \node[above] at (1.5, 2) {\footnotesize 01};
        \node[above] at (2.5, 2) {\footnotesize 11};
        \node[above] at (3.5, 2) {\footnotesize 10};
        \node[left] at (0, 1.5) {\footnotesize 0};
        \node[left] at (0, 0.5) {\footnotesize 1};
        \node[above left] at (0, 2) {\footnotesize #2\textbackslash #3};
        \foreach \val [count=\i from 0] in {#4} {
            \pgfmathtruncatemacro{\col}{mod(\i,4)}
            \pgfmathtruncatemacro{\row}{1 - int(\i/4)}
            \node at (\col+0.5, \row+0.5) {\val};
        }
        #1
    \end{tikzpicture}
}

% 4 değişkenli Karno haritası
% #1: ek TikZ kodu
% #2: satır değişkenleri, #3: sütun değişkenleri
% #4: 16 hücre değeri (AB\CD sıralaması: Karno sırasıyla 00,01,11,10)
\newcommand{\karnodort}[4][]{%
    \begin{tikzpicture}[scale=0.8]
        \draw (0,0) grid (4,4);
        \node[above] at (0.5, 4) {\footnotesize 00};
        \node[above] at (1.5, 4) {\footnotesize 01};
        \node[above] at (2.5, 4) {\footnotesize 11};
        \node[above] at (3.5, 4) {\footnotesize 10};
        \node[left] at (0, 3.5) {\footnotesize 00};
        \node[left] at (0, 2.5) {\footnotesize 01};
        \node[left] at (0, 1.5) {\footnotesize 11};
        \node[left] at (0, 0.5) {\footnotesize 10};
        \node[above left] at (0, 4) {\footnotesize #2\textbackslash #3};
        \foreach \val [count=\i from 0] in {#4} {
            \pgfmathtruncatemacro{\col}{mod(\i,4)}
            \pgfmathtruncatemacro{\row}{3 - int(\i/4)}
            \node at (\col+0.5, \row+0.5) {\val};
        }
        #1
    \end{tikzpicture}
}

% -----------------------------------------------
% FLIP-FLOP MAKROLARI
% -----------------------------------------------

% Genel flip-flop çizim stili
\tikzset{
    flipflop/.style={
        draw, rectangle, minimum width=1.8cm, minimum height=2.4cm,
        font=\small, thick
    },
    ffgiris/.style={font=\footnotesize},
    ffcikis/.style={font=\footnotesize},
}

% -----------------------------------------------
% MANTIK KAPISI STİLLERİ
% -----------------------------------------------

\tikzset{
    % Mantık kapısı sembol boyutları
    kapiboy/.style={
        circuit logic US, huge circuit symbols
    },
}

% -----------------------------------------------
% BIT ALANI DİYAGRAMI MAKROSU
% -----------------------------------------------

% Buyruk bit alanı çizimi
% #1: toplam bit genişliği (örn. 32)
% #2: alanlar listesi (tikz node'ları olarak verilecek)
\tikzset{
    bitalani/.style={
        draw, minimum height=0.8cm, font=\footnotesize\ttfamily,
        anchor=west, fill=blokmavi
    },
    bitust/.style={
        font=\tiny, above, text=gray
    },
    bitalt/.style={
        font=\scriptsize, below
    },
}

% -----------------------------------------------
% AKIŞ ŞEMASI STİLLERİ
% -----------------------------------------------

\tikzset{
    % Akış şeması blokları
    akisbaslangic/.style={
        draw, rounded rectangle, rounded rectangle west arc=10pt,
        rounded rectangle east arc=10pt,
        minimum width=2cm, minimum height=0.6cm,
        fill=blokyesil, font=\footnotesize, align=center,
        thick
    },
    akisislem/.style={
        draw, rectangle, rounded corners=2pt,
        minimum width=2.5cm, minimum height=0.6cm,
        fill=blokmavi, font=\footnotesize, align=center,
        thick
    },
    akiskarar/.style={
        draw, diamond, aspect=2, minimum width=1.8cm,
        fill=bloksari, font=\footnotesize, align=center,
        inner sep=1pt, thick
    },
    akisok/.style={
        ->, >=Stealth, thick
    },
    % Geri dönüş oku (döngü çizgileri için)
    akisdonus/.style={
        ->, >=Stealth, thick, rounded corners=8pt
    },
}

% -----------------------------------------------
% ZAMANLAMA ŞEMASI STİLLERİ
% -----------------------------------------------

\tikzset{
    zamansinyal/.style={
        very thick, sinyalmavi
    },
    zamanok/.style={
        <->, >=Stealth, thin
    },
    zamanetiket/.style={
        font=\footnotesize
    },
}

% -----------------------------------------------
% VERİYOLU BİLEŞEN STİLLERİ (Bölüm 4-5 için)
% -----------------------------------------------

\tikzset{
    % Çoklayıcı (MUX)
    mux/.style={
        draw, trapezium, trapezium angle=70,
        minimum width=0.6cm, minimum height=1.2cm,
        font=\tiny, thick, fill=blokgri, shape border rotate=270
    },
    % AMB
    amb/.style={
        draw, rectangle, minimum width=1.2cm, minimum height=1.8cm,
        fill=blokmavi, font=\small, thick, align=center
    },
    % Toplayıcı (daire)
    toplayici/.style={
        draw, circle, minimum size=0.6cm, font=\small, thick,
        fill=blokgri, inner sep=0pt
    },
    % Yazmaç (ince uzun dikdörtgen)
    yazmac/.style={
        draw, rectangle, minimum width=0.3cm, minimum height=2cm,
        fill=bloksari, font=\tiny, thick
    },
    % Denetim birimi
    denetim/.style={
        draw, ellipse, minimum width=2cm, minimum height=1.2cm,
        fill=blokkirmizi!30, font=\small, thick, align=center
    },
    % Bellek
    bellek/.style={
        draw, rectangle, minimum width=1.8cm, minimum height=2cm,
        fill=blokyesil, font=\small, thick, align=center
    },
    % Veri yolu (kalın çizgi)
    veriyolu/.style={
        line width=2pt, ->, >=Stealth
    },
    % İnce sinyal
    sinyal/.style={
        ->, >=Stealth, thin
    },
}

% -----------------------------------------------
% DURUM MAKİNESİ STİLLERİ
% -----------------------------------------------

\tikzset{
    durum/.style={
        draw, circle, minimum size=1.2cm, font=\small,
        thick, fill=blokmavi
    },
    durumok/.style={
        ->, >=Stealth, thick, bend left=20
    },
    durumdongu/.style={
        ->, >=Stealth, thick, loop above, looseness=8
    },
}

% -----------------------------------------------
% BELLEK DİYAGRAMI STİLLERİ (Bölüm 6 için)
% -----------------------------------------------

\tikzset{
    % Önbellek satırı
    onbellek/.style={
        draw, rectangle, minimum height=0.5cm, font=\tiny,
        fill=blokmavi!30
    },
    % Bellek hücresi
    bellekhucre/.style={
        draw, rectangle, minimum width=0.8cm, minimum height=0.5cm,
        font=\tiny
    },
}

% -----------------------------------------------
% BORU HATTI DİYAGRAM STİLLERİ (Bölüm 5-6 için)
% -----------------------------------------------

\tikzset{
    % Pipeline aşama kutusu
    asama/.style={
        draw, rectangle, minimum width=1.2cm, minimum height=0.7cm,
        font=\footnotesize, align=center, thick
    },
    % Beş aşama stili
    asamaG/.style={asama, fill=asamagetir},
    asamaC/.style={asama, fill=asamacoz},
    asamaU/.style={asama, fill=asamayurut},
    asamaB/.style={asama, fill=asamabellek},
    asamaY/.style={asama, fill=asamayaz},
    % Boş/stall aşama
    asamabos/.style={asama, fill=blokgri, dashed},
    % Pipeline yazmaç (aşamalar arası)
    pipeyazmac/.style={
        draw, rectangle, minimum width=0.15cm, minimum height=1.8cm,
        fill=bloksari, thick
    },
    % İşlemci birimi (veriyolu bileşeni)
    birim/.style={
        draw, rectangle, rounded corners=3pt,
        minimum width=1.5cm, minimum height=1cm,
        fill=blokmavi, font=\small, thick, align=center
    },
}

% -----------------------------------------------
% ÇOK ÇEKİRDEKLİ / PARALEL STİLLER (Bölüm 8-9)
% -----------------------------------------------

\tikzset{
    % İşlemci çekirdeği
    cekirdek/.style={
        draw, rectangle, rounded corners=4pt,
        minimum width=2cm, minimum height=1.5cm,
        fill=blokmavi, font=\small, thick, align=center
    },
    % Önbellek kutusu
    onbellekblok/.style={
        draw, rectangle, minimum width=1.8cm, minimum height=0.8cm,
        fill=blokyesil-acik, font=\footnotesize, align=center
    },
    % Paylaşılan kaynak
    paylasilan/.style={
        draw, rectangle, rounded corners=2pt,
        minimum width=3cm, minimum height=1cm,
        fill=bloksari, font=\small, thick, align=center
    },
    % Bağlantı yolu (bus/interconnect)
    baglanti/.style={
        draw, rectangle, minimum height=0.4cm,
        fill=blokgri, font=\footnotesize
    },
}

% -----------------------------------------------
% GPU / HIZLANDIRICI STİLLERİ (Bölüm 10)
% -----------------------------------------------

\tikzset{
    % SM / CU birimi
    gpubirim/.style={
        draw, rectangle, rounded corners=2pt,
        minimum width=1cm, minimum height=0.8cm,
        fill=blokmavi, font=\tiny, thick, align=center
    },
    % Bellek modülü
    gpubellek/.style={
        draw, rectangle, minimum width=2cm, minimum height=0.6cm,
        fill=blokyesil, font=\footnotesize, align=center
    },
}
